\documentclass[authoryear,1p,12pt]{elsarticle}
\usepackage{amsmath}
\usepackage{amssymb}
\usepackage[mathlines,displaymath]{lineno}
\usepackage{graphicx}
\usepackage[all,2cell,dvips]{xy}
\usepackage{natbib}
\usepackage[usenames,rgb]{xcolor} 
\usepackage{ucs}                                            
\usepackage[utf8]{inputenc}                                                                           
\usepackage{array}                                            
\usepackage{longtable}                                        
\usepackage{calc}                                             
\usepackage{multirow}                                         
\usepackage{hhline}                                           
\usepackage{ifthen} 
% \usepackage{babel}
\usepackage{tcolorbox}
\newtcolorbox{mybox}{colback=grey!5!white,colframe=red!75!black}
\usepackage{setspace}
\usepackage{verbatim}
\usepackage{etaremune}
\usepackage{url}
\usepackage{color}
\usepackage{graphicx}
\usepackage{natbib}
\usepackage[latin1]{inputenc}
\usepackage{color}
\usepackage{array}
\usepackage{float}
\usepackage{wrapfig}

\newcommand{\etal}{{et~al.{}}}
\newcommand{\ie}{{i.~e.{}}}
\newcommand{\eg}{{e.~g.{}}}
\newcommand{\viz}{{viz.{}}}
\newcommand{\etc}{{etc.{}}}
\newcommand{\apriori}{{a priori{}}}
\newcommand{\vv}{{vice versa{}}}
\newcommand{\cf}{cf.{}}
\setcounter{tocdepth}{3}
\usepackage{tikz}
\newcommand{\carlos}[1]{\textcolor{Red}{#1}}

% Text layout
\topmargin 0.0cm
\oddsidemargin 0.2cm
\evensidemargin 0.2cm
\textwidth 16cm 
\textheight 20.5cm


\begin{document}

%\begin{mybox}\begin{singlespace}
%\begin{small}
\section*{Box 1: Metaecosystem multitrait biodiversity model}

We outline a multitrait-based model to evaluate the effect of trait
architecture, i.e., the covariation among biotic, abiotic and
migration traits on biodiversity rescue (Figure 1). We explore two
trait architectures rooted in empirical and theoretical studies: the
modular and the magic trait integration architecture
\citep{ENDLER199522,Armbrusteretal2018}. The modular scenario takes
into account independent traits, i.e. each trait contributes
independently to fitness (Figure 1, Modular). The magic trait
architecture incorporates the biotic trait as the main contributor to
fitness with the abiotic and the migration traits depending on the
evolution of the biotic trait (Fig 1, Magic).

Our model considers a landscape with a network of patches connected by
migration events and inhabited by resource and consumer populations
comprising several individuals each. Biotic, abiotic and migration
trait values are initially sampled from normal distributions each with
mean and variance for each trait given by
$\mathcal{N}(\mu_{b},\sigma_{b}^2)$,
$\mathcal{N}(\mu_{a},\sigma_{a}^2)$, and
$\mathcal{N}(\mu_{m},\sigma_{m}^2)$, for the biotic, abiotic, and
migration trait, respectively. Population dynamics of resource species
$i$ is driven by the trait- and temporal-dependent fitness functions
following \vspace{0.25 in}
  \begin{equation}
    \hspace{0.9 in} \mathcal{W}_{R}({\bf z_{B_{i}}})^{t}_{x} = 1 /(1 + exp(\left -\gamma(\mathrm{\bf z_{B_{i}}}^{t}_{x} - {\bf \hat{z}_{B_{for all_{j}}}}^{t}_{x})^2)), \nonumber
  \end{equation}
  \begin{equation}
    \mathcal{W}_{R}({\bf z_{A_{i}}})^{t}_{x} = exp[- \gamma(\mathrm{\bf z_{A_{i}}}^{t}_{x} - \mathrm{\bf \hat{z}_{A_{i}}}_{x})^2], %\nonumber
  \end{equation}
  \begin{equation}
    \mathcal{W}_{R}({\bf z_{M_{i}}})^{t}_{x} = exp[- \gamma(\mathrm{\bf z_{M_{i}}}^{t}_{x} - \mathrm{\bf \hat{\theta}_{D}}_{x})^2]. \nonumber
    \end{equation}
    \vspace{0.25 in}
    
    Similarly, the population dynamics for the consumer species is
    given by
    
\vspace{0.25 in}
  \begin{equation}
   \hspace{0.25 in} \mathcal{W}_{C}({\bf z_{B_{i}}})^{t}_{x} = exp[-\gamma(\mathrm{\bf z_{B_{i}}}^{t}_{x} - {\bf \hat{z}_{B_{for all_{j}}}}^{t}_{x})^2], \nonumber
  \end{equation}
\begin{equation}
  \mathcal{W}_{C}({\bf z_{A_{i}}})^{t}_{x} = exp[- \gamma(\mathrm{\bf z_{A_{i}}}^{t}_{x} - \mathrm{\bf \hat{z}_{A_{i}}}_{x})^2], %\nonumber
\end{equation}
\begin{equation}
  \mathcal{W}_{C}({\bf z_{M_{i}}})^{t}_{x} = exp[- \gamma(\mathrm{\bf z_{M_{i}}}^{t}_{x} - \mathrm{\bf \hat{\theta}_{D}}_{x})^2], \nonumber
  \end{equation}

  \vspace{0.25 in} where $\mathrm{\bf z_{B_{i}}}^{t}_{x}$,
  $\mathrm{\bf z_{A_{i}}}^{t}_{x}$, and
  $\mathrm{\bf z_{M_{i}}}^{t}_{x}$ represent the biotic, abiotic, and
  migration trait value of phenotype $i$ at time $t$ and site $x$ for
  resources and consumer species, $R$ and $C$, respectively. Equation
  1 can be compactly written for resources as
  
  % Resources
  \begin{center}
\begin{equation}
  \hspace{0.05 in}\mathcal{W}_{R}({\bf z_{BAM_{i}}})^{t}_{x} = \mathcal{E}[-\gamma([{\bf z_{\alpha_{i}}}^{t}_{x} - {\bf \theta^{*}}_{x}]^T} {\bf \Omega_{BAM}}^{-1} [{\bf z_{\alpha_{i}}}^{t}_{x} - {\bf \theta^{*}}_{x}]})], \hspace{1 in}
\end{equation}
\end{center}
\vspace{0.25 in}
and for consumers as
  \begin{center}
\begin{equation}
  \hspace{0.05 in}\mathcal{W}_{C}({\bf z_{BAM_{i}}})^{t}_{x} = exp[-\gamma([{\bf z_{\alpha_{i}}}^{t}_{x} - {\bf \theta^{*}}_{x}]^T} {\bf \Omega_{BAM}}^{-1} [{\bf z_{\alpha_{i}}}^{t}_{x} - {\bf \theta^{*}}_{x}]})], \hspace{1 in}
\end{equation}
\end{center}
\vspace{0.2 in} where $\mathcal{E}$ in the resources equation
represents the fact that the biotic trait value increases its fitness
the more it deviates from the mean and common predator type (i.e.,
negative-frequency dependent selection, equation 1).
${\bf \theta^{*}}_{x}$ is the multivariate fitness optimum for the
biotic, abiotic and migration trait in site $x$, with the asterisk
referring to being time-dependent in the biotic, but not in the
abiotic and migration traits. ${\bf \Omega}$ is the covariance matrix
\citep{Lande:1980, Melo&Marroig:2014}, and $\gamma$ determines the
interaction sensitivity to deviations from the biotic, abiotic and
migration optimum. The interpretation of $\gamma$ is as follows: it
represents the strength of coevolutionary selection between resource
and consumer traits (i.e., $\gamma$, y-axis in Figure 1), stabilizing
environmental selection and stabilizing migration selection for the
biotic, abiotic and migration traits, respectively.

If the covariance matrix, ${\bf \Omega}$, is diagonal, representing
our modular scenario (Fig. 1, left), then there is no correlated
stabilizing selection (i.e., Modular scenario was run with
off-diagonal values equal to zero, ${\bf \Omega}$ = [1 0 0;0 1 0;0 0
1];). If the covariance matrix is not diagonal, our magic scenario
(Fig. 1, right), then there is correlated stabilizing selection with
the biotic trait determining the selection on the other traits (i.e.,
Magic scenario was run with the following covariance matrix,
${\bf \Omega}$ = [1 0.7 0.7;0.7 1 0;0.7 0 1]). The biotic optimum is
time-dependent and we define it as the maximum and minimum mean trait
distance across all populations in site $x$ for the resource and
consumer biotic trait value, respectively (Biotic trait, equation 1
and 2). Migration optimum is given by the mean distance in the
landscape (i.e., each site has coordinates from which we calculate the
Euclidian distance between each pair of sites). Phenotypes with
migration trait value near to this optimum maximize fitness for this
trait only for the modular but not for the magic trait scenario. The
ecology underpinning this migration optimum takes into account that
non migrating or migrating excesively is more risky than the optimal
migration. The abiotic optimum is taken initially as the mean of the
distribution for each species and it remains the same during the
dynamics.
  \begin{center}
  \begin{figure}[h!]
   \hspace{-0.15 in}\includegraphics[width=1.05\textwidth]{Heatmap.pdf}
  \vspace{0.1 in}
  \caption{\small {\bf Effect of modular and magic trait architecture
      on biodiversity rescue}: We evaluate the importance of
    coevolutionary selection and migration as biodiversity rescue
    mechanisms in a Modular (left) and Magic (right), trait
    architecture scenarios. We explore the formation of hot and cold
    spots of species persistence in the landscape along a gradient of
    migration and coevolutionary selection strength. The modular
    scenario predicts higher biodiversity rescue than the magic one
    along the migration and coevolutionary selection gradient. Hot
    spots in biodiversity rescue occur for low-to-medium migration and
    medium-to-high coevolutionary selection values. The magic scenario
    predicts a highly homogeneous biodiversity rescue along the
    gradients}
\end{figure}
\end{center}
\vspace{-0.5 in} Population dynamics run as follows: at each time step
the time-dependent fitness function in site $x$, for resource species
$i$, $\mathcal{W}_{R}({\bf z_{BAM_{i}}})}^{t}_{x}$, and consumer
species $j$, $\mathcal{W}_{C}({\bf z_{BAM_{j}}})}^{t}_{x}$, determines
the probabilities to die. Individuals that die are replaced by
offspring from the same site $x$ or from migration with probability
${\bf m}$. Offspring's phenotypes, {\bf $z_o$} =
$[z_{ob}, z_{oa}, z_{om}]$, are inherited from the mother plus an
environmental deviation, {\bf e}, taken from a multivariate uniform
distribution with range $\pm0.001$ in all dimensions (i.e., these
values are independently sampled for the modular but correlated
following the input covariance matrix for the magic trait
scenario). We run our model for 500 replicates each for 50 generations
for the two scenarios. To evaluate the importance of coevolution and
migration as biodiversity rescue mechanisms, we plot the formation of
hot and cold spots of species frequency along a gradient of migration
and coevolutionary selection strength. Simulations were run for 10
sites in the landscape each containing 1000 individuals with equal
initial abundances for 5 resource and 5 consumer species. Initial mean
and variance for the biotic, abiotic, and migration traits for
resources and consumers were drawn from normal distributions,
$\mathcal{N}(2,12)$, $\mathcal{N}(50,9)$, and
$\mathcal{N}(60,2)$. Codes and outputs are available from
https://github.com/melian009/Cometa).


\begin{comment}
Mean species richness decays with the strength of coevolutionary
selection for the biotic trait in the two scenarios, the modular
(black) and the integrated (red) scenario (Figure 1b). Adding
interactions among traits, the integrated scenario, decreases the
number of species for all the gradient of coevolutionary selection,
compared to the modular scenario.

Q2: mention the initial values: number of species, population size,
initial trait variance and covariance values, and the number of
generations per replicate.


Armbruster This study demonstrates the value of allometric and
variance analyses (in addition to correlation) in assessing the
evolutionary significance of floral-trait stability and plasticity.
\end{comment}

% \end{small}
%\end{singlespace}
%\end{mybox}

\newpage
\section{Bibliography}
\bibliographystyle{unsrtnat}
\bibliography{space.bib}
\end{document}